\documentclass[11pt]{article}

    \usepackage[breakable]{tcolorbox}
    \usepackage{parskip} % Stop auto-indenting (to mimic markdown behaviour)
    

    % Basic figure setup, for now with no caption control since it's done
    % automatically by Pandoc (which extracts ![](path) syntax from Markdown).
    \usepackage{graphicx}
    % Keep aspect ratio if custom image width or height is specified
    \setkeys{Gin}{keepaspectratio}
    % Maintain compatibility with old templates. Remove in nbconvert 6.0
    \let\Oldincludegraphics\includegraphics
    % Ensure that by default, figures have no caption (until we provide a
    % proper Figure object with a Caption API and a way to capture that
    % in the conversion process - todo).
    \usepackage{caption}
    \DeclareCaptionFormat{nocaption}{}
    \captionsetup{format=nocaption,aboveskip=0pt,belowskip=0pt}

    \usepackage{float}
    \floatplacement{figure}{H} % forces figures to be placed at the correct location
    \usepackage{xcolor} % Allow colors to be defined
    \usepackage{enumerate} % Needed for markdown enumerations to work
    \usepackage{geometry} % Used to adjust the document margins
    \usepackage{amsmath} % Equations
    \usepackage{amssymb} % Equations
    \usepackage{textcomp} % defines textquotesingle
    % Hack from http://tex.stackexchange.com/a/47451/13684:
    \AtBeginDocument{%
        \def\PYZsq{\textquotesingle}% Upright quotes in Pygmentized code
    }
    \usepackage{upquote} % Upright quotes for verbatim code
    \usepackage{eurosym} % defines \euro

    \usepackage{iftex}
    \ifPDFTeX
        \usepackage[T1]{fontenc}
        \IfFileExists{alphabeta.sty}{
              \usepackage{alphabeta}
          }{
              \usepackage[mathletters]{ucs}
              \usepackage[utf8x]{inputenc}
          }
    \else
        \usepackage{fontspec}
        \usepackage{unicode-math}
    \fi

    \usepackage{fancyvrb} % verbatim replacement that allows latex
    \usepackage{grffile} % extends the file name processing of package graphics
                         % to support a larger range
    \makeatletter % fix for old versions of grffile with XeLaTeX
    \@ifpackagelater{grffile}{2019/11/01}
    {
      % Do nothing on new versions
    }
    {
      \def\Gread@@xetex#1{%
        \IfFileExists{"\Gin@base".bb}%
        {\Gread@eps{\Gin@base.bb}}%
        {\Gread@@xetex@aux#1}%
      }
    }
    \makeatother
    \usepackage[Export]{adjustbox} % Used to constrain images to a maximum size
    \adjustboxset{max size={0.9\linewidth}{0.9\paperheight}}

    % The hyperref package gives us a pdf with properly built
    % internal navigation ('pdf bookmarks' for the table of contents,
    % internal cross-reference links, web links for URLs, etc.)
    \usepackage{hyperref}
    % The default LaTeX title has an obnoxious amount of whitespace. By default,
    % titling removes some of it. It also provides customization options.
    \usepackage{titling}
    \usepackage{longtable} % longtable support required by pandoc >1.10
    \usepackage{booktabs}  % table support for pandoc > 1.12.2
    \usepackage{array}     % table support for pandoc >= 2.11.3
    \usepackage{calc}      % table minipage width calculation for pandoc >= 2.11.1
    \usepackage[inline]{enumitem} % IRkernel/repr support (it uses the enumerate* environment)
    \usepackage[normalem]{ulem} % ulem is needed to support strikethroughs (\sout)
                                % normalem makes italics be italics, not underlines
    \usepackage{soul}      % strikethrough (\st) support for pandoc >= 3.0.0
    \usepackage{mathrsfs}
    

    
    % Colors for the hyperref package
    \definecolor{urlcolor}{rgb}{0,.145,.698}
    \definecolor{linkcolor}{rgb}{.71,0.21,0.01}
    \definecolor{citecolor}{rgb}{.12,.54,.11}

    % ANSI colors
    \definecolor{ansi-black}{HTML}{3E424D}
    \definecolor{ansi-black-intense}{HTML}{282C36}
    \definecolor{ansi-red}{HTML}{E75C58}
    \definecolor{ansi-red-intense}{HTML}{B22B31}
    \definecolor{ansi-green}{HTML}{00A250}
    \definecolor{ansi-green-intense}{HTML}{007427}
    \definecolor{ansi-yellow}{HTML}{DDB62B}
    \definecolor{ansi-yellow-intense}{HTML}{B27D12}
    \definecolor{ansi-blue}{HTML}{208FFB}
    \definecolor{ansi-blue-intense}{HTML}{0065CA}
    \definecolor{ansi-magenta}{HTML}{D160C4}
    \definecolor{ansi-magenta-intense}{HTML}{A03196}
    \definecolor{ansi-cyan}{HTML}{60C6C8}
    \definecolor{ansi-cyan-intense}{HTML}{258F8F}
    \definecolor{ansi-white}{HTML}{C5C1B4}
    \definecolor{ansi-white-intense}{HTML}{A1A6B2}
    \definecolor{ansi-default-inverse-fg}{HTML}{FFFFFF}
    \definecolor{ansi-default-inverse-bg}{HTML}{000000}

    % common color for the border for error outputs.
    \definecolor{outerrorbackground}{HTML}{FFDFDF}

    % commands and environments needed by pandoc snippets
    % extracted from the output of `pandoc -s`
    \providecommand{\tightlist}{%
      \setlength{\itemsep}{0pt}\setlength{\parskip}{0pt}}
    \DefineVerbatimEnvironment{Highlighting}{Verbatim}{commandchars=\\\{\}}
    % Add ',fontsize=\small' for more characters per line
    \newenvironment{Shaded}{}{}
    \newcommand{\KeywordTok}[1]{\textcolor[rgb]{0.00,0.44,0.13}{\textbf{{#1}}}}
    \newcommand{\DataTypeTok}[1]{\textcolor[rgb]{0.56,0.13,0.00}{{#1}}}
    \newcommand{\DecValTok}[1]{\textcolor[rgb]{0.25,0.63,0.44}{{#1}}}
    \newcommand{\BaseNTok}[1]{\textcolor[rgb]{0.25,0.63,0.44}{{#1}}}
    \newcommand{\FloatTok}[1]{\textcolor[rgb]{0.25,0.63,0.44}{{#1}}}
    \newcommand{\CharTok}[1]{\textcolor[rgb]{0.25,0.44,0.63}{{#1}}}
    \newcommand{\StringTok}[1]{\textcolor[rgb]{0.25,0.44,0.63}{{#1}}}
    \newcommand{\CommentTok}[1]{\textcolor[rgb]{0.38,0.63,0.69}{\textit{{#1}}}}
    \newcommand{\OtherTok}[1]{\textcolor[rgb]{0.00,0.44,0.13}{{#1}}}
    \newcommand{\AlertTok}[1]{\textcolor[rgb]{1.00,0.00,0.00}{\textbf{{#1}}}}
    \newcommand{\FunctionTok}[1]{\textcolor[rgb]{0.02,0.16,0.49}{{#1}}}
    \newcommand{\RegionMarkerTok}[1]{{#1}}
    \newcommand{\ErrorTok}[1]{\textcolor[rgb]{1.00,0.00,0.00}{\textbf{{#1}}}}
    \newcommand{\NormalTok}[1]{{#1}}

    % Additional commands for more recent versions of Pandoc
    \newcommand{\ConstantTok}[1]{\textcolor[rgb]{0.53,0.00,0.00}{{#1}}}
    \newcommand{\SpecialCharTok}[1]{\textcolor[rgb]{0.25,0.44,0.63}{{#1}}}
    \newcommand{\VerbatimStringTok}[1]{\textcolor[rgb]{0.25,0.44,0.63}{{#1}}}
    \newcommand{\SpecialStringTok}[1]{\textcolor[rgb]{0.73,0.40,0.53}{{#1}}}
    \newcommand{\ImportTok}[1]{{#1}}
    \newcommand{\DocumentationTok}[1]{\textcolor[rgb]{0.73,0.13,0.13}{\textit{{#1}}}}
    \newcommand{\AnnotationTok}[1]{\textcolor[rgb]{0.38,0.63,0.69}{\textbf{\textit{{#1}}}}}
    \newcommand{\CommentVarTok}[1]{\textcolor[rgb]{0.38,0.63,0.69}{\textbf{\textit{{#1}}}}}
    \newcommand{\VariableTok}[1]{\textcolor[rgb]{0.10,0.09,0.49}{{#1}}}
    \newcommand{\ControlFlowTok}[1]{\textcolor[rgb]{0.00,0.44,0.13}{\textbf{{#1}}}}
    \newcommand{\OperatorTok}[1]{\textcolor[rgb]{0.40,0.40,0.40}{{#1}}}
    \newcommand{\BuiltInTok}[1]{{#1}}
    \newcommand{\ExtensionTok}[1]{{#1}}
    \newcommand{\PreprocessorTok}[1]{\textcolor[rgb]{0.74,0.48,0.00}{{#1}}}
    \newcommand{\AttributeTok}[1]{\textcolor[rgb]{0.49,0.56,0.16}{{#1}}}
    \newcommand{\InformationTok}[1]{\textcolor[rgb]{0.38,0.63,0.69}{\textbf{\textit{{#1}}}}}
    \newcommand{\WarningTok}[1]{\textcolor[rgb]{0.38,0.63,0.69}{\textbf{\textit{{#1}}}}}
    \makeatletter
    \newsavebox\pandoc@box
    \newcommand*\pandocbounded[1]{%
      \sbox\pandoc@box{#1}%
      % scaling factors for width and height
      \Gscale@div\@tempa\textheight{\dimexpr\ht\pandoc@box+\dp\pandoc@box\relax}%
      \Gscale@div\@tempb\linewidth{\wd\pandoc@box}%
      % select the smaller of both
      \ifdim\@tempb\p@<\@tempa\p@
        \let\@tempa\@tempb
      \fi
      % scaling accordingly (\@tempa < 1)
      \ifdim\@tempa\p@<\p@
        \scalebox{\@tempa}{\usebox\pandoc@box}%
      % scaling not needed, use as it is
      \else
        \usebox{\pandoc@box}%
      \fi
    }
    \makeatother

    % Define a nice break command that doesn't care if a line doesn't already
    % exist.
    \def\br{\hspace*{\fill} \\* }
    % Math Jax compatibility definitions
    \def\gt{>}
    \def\lt{<}
    \let\Oldtex\TeX
    \let\Oldlatex\LaTeX
    \renewcommand{\TeX}{\textrm{\Oldtex}}
    \renewcommand{\LaTeX}{\textrm{\Oldlatex}}
    % Document parameters
    % Document title
    \title{01-mandat}
    
    
    
    
    
    
    
% Pygments definitions
\makeatletter
\def\PY@reset{\let\PY@it=\relax \let\PY@bf=\relax%
    \let\PY@ul=\relax \let\PY@tc=\relax%
    \let\PY@bc=\relax \let\PY@ff=\relax}
\def\PY@tok#1{\csname PY@tok@#1\endcsname}
\def\PY@toks#1+{\ifx\relax#1\empty\else%
    \PY@tok{#1}\expandafter\PY@toks\fi}
\def\PY@do#1{\PY@bc{\PY@tc{\PY@ul{%
    \PY@it{\PY@bf{\PY@ff{#1}}}}}}}
\def\PY#1#2{\PY@reset\PY@toks#1+\relax+\PY@do{#2}}

\@namedef{PY@tok@w}{\def\PY@tc##1{\textcolor[rgb]{0.73,0.73,0.73}{##1}}}
\@namedef{PY@tok@c}{\let\PY@it=\textit\def\PY@tc##1{\textcolor[rgb]{0.24,0.48,0.48}{##1}}}
\@namedef{PY@tok@cp}{\def\PY@tc##1{\textcolor[rgb]{0.61,0.40,0.00}{##1}}}
\@namedef{PY@tok@k}{\let\PY@bf=\textbf\def\PY@tc##1{\textcolor[rgb]{0.00,0.50,0.00}{##1}}}
\@namedef{PY@tok@kp}{\def\PY@tc##1{\textcolor[rgb]{0.00,0.50,0.00}{##1}}}
\@namedef{PY@tok@kt}{\def\PY@tc##1{\textcolor[rgb]{0.69,0.00,0.25}{##1}}}
\@namedef{PY@tok@o}{\def\PY@tc##1{\textcolor[rgb]{0.40,0.40,0.40}{##1}}}
\@namedef{PY@tok@ow}{\let\PY@bf=\textbf\def\PY@tc##1{\textcolor[rgb]{0.67,0.13,1.00}{##1}}}
\@namedef{PY@tok@nb}{\def\PY@tc##1{\textcolor[rgb]{0.00,0.50,0.00}{##1}}}
\@namedef{PY@tok@nf}{\def\PY@tc##1{\textcolor[rgb]{0.00,0.00,1.00}{##1}}}
\@namedef{PY@tok@nc}{\let\PY@bf=\textbf\def\PY@tc##1{\textcolor[rgb]{0.00,0.00,1.00}{##1}}}
\@namedef{PY@tok@nn}{\let\PY@bf=\textbf\def\PY@tc##1{\textcolor[rgb]{0.00,0.00,1.00}{##1}}}
\@namedef{PY@tok@ne}{\let\PY@bf=\textbf\def\PY@tc##1{\textcolor[rgb]{0.80,0.25,0.22}{##1}}}
\@namedef{PY@tok@nv}{\def\PY@tc##1{\textcolor[rgb]{0.10,0.09,0.49}{##1}}}
\@namedef{PY@tok@no}{\def\PY@tc##1{\textcolor[rgb]{0.53,0.00,0.00}{##1}}}
\@namedef{PY@tok@nl}{\def\PY@tc##1{\textcolor[rgb]{0.46,0.46,0.00}{##1}}}
\@namedef{PY@tok@ni}{\let\PY@bf=\textbf\def\PY@tc##1{\textcolor[rgb]{0.44,0.44,0.44}{##1}}}
\@namedef{PY@tok@na}{\def\PY@tc##1{\textcolor[rgb]{0.41,0.47,0.13}{##1}}}
\@namedef{PY@tok@nt}{\let\PY@bf=\textbf\def\PY@tc##1{\textcolor[rgb]{0.00,0.50,0.00}{##1}}}
\@namedef{PY@tok@nd}{\def\PY@tc##1{\textcolor[rgb]{0.67,0.13,1.00}{##1}}}
\@namedef{PY@tok@s}{\def\PY@tc##1{\textcolor[rgb]{0.73,0.13,0.13}{##1}}}
\@namedef{PY@tok@sd}{\let\PY@it=\textit\def\PY@tc##1{\textcolor[rgb]{0.73,0.13,0.13}{##1}}}
\@namedef{PY@tok@si}{\let\PY@bf=\textbf\def\PY@tc##1{\textcolor[rgb]{0.64,0.35,0.47}{##1}}}
\@namedef{PY@tok@se}{\let\PY@bf=\textbf\def\PY@tc##1{\textcolor[rgb]{0.67,0.36,0.12}{##1}}}
\@namedef{PY@tok@sr}{\def\PY@tc##1{\textcolor[rgb]{0.64,0.35,0.47}{##1}}}
\@namedef{PY@tok@ss}{\def\PY@tc##1{\textcolor[rgb]{0.10,0.09,0.49}{##1}}}
\@namedef{PY@tok@sx}{\def\PY@tc##1{\textcolor[rgb]{0.00,0.50,0.00}{##1}}}
\@namedef{PY@tok@m}{\def\PY@tc##1{\textcolor[rgb]{0.40,0.40,0.40}{##1}}}
\@namedef{PY@tok@gh}{\let\PY@bf=\textbf\def\PY@tc##1{\textcolor[rgb]{0.00,0.00,0.50}{##1}}}
\@namedef{PY@tok@gu}{\let\PY@bf=\textbf\def\PY@tc##1{\textcolor[rgb]{0.50,0.00,0.50}{##1}}}
\@namedef{PY@tok@gd}{\def\PY@tc##1{\textcolor[rgb]{0.63,0.00,0.00}{##1}}}
\@namedef{PY@tok@gi}{\def\PY@tc##1{\textcolor[rgb]{0.00,0.52,0.00}{##1}}}
\@namedef{PY@tok@gr}{\def\PY@tc##1{\textcolor[rgb]{0.89,0.00,0.00}{##1}}}
\@namedef{PY@tok@ge}{\let\PY@it=\textit}
\@namedef{PY@tok@gs}{\let\PY@bf=\textbf}
\@namedef{PY@tok@ges}{\let\PY@bf=\textbf\let\PY@it=\textit}
\@namedef{PY@tok@gp}{\let\PY@bf=\textbf\def\PY@tc##1{\textcolor[rgb]{0.00,0.00,0.50}{##1}}}
\@namedef{PY@tok@go}{\def\PY@tc##1{\textcolor[rgb]{0.44,0.44,0.44}{##1}}}
\@namedef{PY@tok@gt}{\def\PY@tc##1{\textcolor[rgb]{0.00,0.27,0.87}{##1}}}
\@namedef{PY@tok@err}{\def\PY@bc##1{{\setlength{\fboxsep}{\string -\fboxrule}\fcolorbox[rgb]{1.00,0.00,0.00}{1,1,1}{\strut ##1}}}}
\@namedef{PY@tok@kc}{\let\PY@bf=\textbf\def\PY@tc##1{\textcolor[rgb]{0.00,0.50,0.00}{##1}}}
\@namedef{PY@tok@kd}{\let\PY@bf=\textbf\def\PY@tc##1{\textcolor[rgb]{0.00,0.50,0.00}{##1}}}
\@namedef{PY@tok@kn}{\let\PY@bf=\textbf\def\PY@tc##1{\textcolor[rgb]{0.00,0.50,0.00}{##1}}}
\@namedef{PY@tok@kr}{\let\PY@bf=\textbf\def\PY@tc##1{\textcolor[rgb]{0.00,0.50,0.00}{##1}}}
\@namedef{PY@tok@bp}{\def\PY@tc##1{\textcolor[rgb]{0.00,0.50,0.00}{##1}}}
\@namedef{PY@tok@fm}{\def\PY@tc##1{\textcolor[rgb]{0.00,0.00,1.00}{##1}}}
\@namedef{PY@tok@vc}{\def\PY@tc##1{\textcolor[rgb]{0.10,0.09,0.49}{##1}}}
\@namedef{PY@tok@vg}{\def\PY@tc##1{\textcolor[rgb]{0.10,0.09,0.49}{##1}}}
\@namedef{PY@tok@vi}{\def\PY@tc##1{\textcolor[rgb]{0.10,0.09,0.49}{##1}}}
\@namedef{PY@tok@vm}{\def\PY@tc##1{\textcolor[rgb]{0.10,0.09,0.49}{##1}}}
\@namedef{PY@tok@sa}{\def\PY@tc##1{\textcolor[rgb]{0.73,0.13,0.13}{##1}}}
\@namedef{PY@tok@sb}{\def\PY@tc##1{\textcolor[rgb]{0.73,0.13,0.13}{##1}}}
\@namedef{PY@tok@sc}{\def\PY@tc##1{\textcolor[rgb]{0.73,0.13,0.13}{##1}}}
\@namedef{PY@tok@dl}{\def\PY@tc##1{\textcolor[rgb]{0.73,0.13,0.13}{##1}}}
\@namedef{PY@tok@s2}{\def\PY@tc##1{\textcolor[rgb]{0.73,0.13,0.13}{##1}}}
\@namedef{PY@tok@sh}{\def\PY@tc##1{\textcolor[rgb]{0.73,0.13,0.13}{##1}}}
\@namedef{PY@tok@s1}{\def\PY@tc##1{\textcolor[rgb]{0.73,0.13,0.13}{##1}}}
\@namedef{PY@tok@mb}{\def\PY@tc##1{\textcolor[rgb]{0.40,0.40,0.40}{##1}}}
\@namedef{PY@tok@mf}{\def\PY@tc##1{\textcolor[rgb]{0.40,0.40,0.40}{##1}}}
\@namedef{PY@tok@mh}{\def\PY@tc##1{\textcolor[rgb]{0.40,0.40,0.40}{##1}}}
\@namedef{PY@tok@mi}{\def\PY@tc##1{\textcolor[rgb]{0.40,0.40,0.40}{##1}}}
\@namedef{PY@tok@il}{\def\PY@tc##1{\textcolor[rgb]{0.40,0.40,0.40}{##1}}}
\@namedef{PY@tok@mo}{\def\PY@tc##1{\textcolor[rgb]{0.40,0.40,0.40}{##1}}}
\@namedef{PY@tok@ch}{\let\PY@it=\textit\def\PY@tc##1{\textcolor[rgb]{0.24,0.48,0.48}{##1}}}
\@namedef{PY@tok@cm}{\let\PY@it=\textit\def\PY@tc##1{\textcolor[rgb]{0.24,0.48,0.48}{##1}}}
\@namedef{PY@tok@cpf}{\let\PY@it=\textit\def\PY@tc##1{\textcolor[rgb]{0.24,0.48,0.48}{##1}}}
\@namedef{PY@tok@c1}{\let\PY@it=\textit\def\PY@tc##1{\textcolor[rgb]{0.24,0.48,0.48}{##1}}}
\@namedef{PY@tok@cs}{\let\PY@it=\textit\def\PY@tc##1{\textcolor[rgb]{0.24,0.48,0.48}{##1}}}

\def\PYZbs{\char`\\}
\def\PYZus{\char`\_}
\def\PYZob{\char`\{}
\def\PYZcb{\char`\}}
\def\PYZca{\char`\^}
\def\PYZam{\char`\&}
\def\PYZlt{\char`\<}
\def\PYZgt{\char`\>}
\def\PYZsh{\char`\#}
\def\PYZpc{\char`\%}
\def\PYZdl{\char`\$}
\def\PYZhy{\char`\-}
\def\PYZsq{\char`\'}
\def\PYZdq{\char`\"}
\def\PYZti{\char`\~}
% for compatibility with earlier versions
\def\PYZat{@}
\def\PYZlb{[}
\def\PYZrb{]}
\makeatother


    % For linebreaks inside Verbatim environment from package fancyvrb.
    \makeatletter
        \newbox\Wrappedcontinuationbox
        \newbox\Wrappedvisiblespacebox
        \newcommand*\Wrappedvisiblespace {\textcolor{red}{\textvisiblespace}}
        \newcommand*\Wrappedcontinuationsymbol {\textcolor{red}{\llap{\tiny$\m@th\hookrightarrow$}}}
        \newcommand*\Wrappedcontinuationindent {3ex }
        \newcommand*\Wrappedafterbreak {\kern\Wrappedcontinuationindent\copy\Wrappedcontinuationbox}
        % Take advantage of the already applied Pygments mark-up to insert
        % potential linebreaks for TeX processing.
        %        {, <, #, %, $, ' and ": go to next line.
        %        _, }, ^, &, >, - and ~: stay at end of broken line.
        % Use of \textquotesingle for straight quote.
        \newcommand*\Wrappedbreaksatspecials {%
            \def\PYGZus{\discretionary{\char`\_}{\Wrappedafterbreak}{\char`\_}}%
            \def\PYGZob{\discretionary{}{\Wrappedafterbreak\char`\{}{\char`\{}}%
            \def\PYGZcb{\discretionary{\char`\}}{\Wrappedafterbreak}{\char`\}}}%
            \def\PYGZca{\discretionary{\char`\^}{\Wrappedafterbreak}{\char`\^}}%
            \def\PYGZam{\discretionary{\char`\&}{\Wrappedafterbreak}{\char`\&}}%
            \def\PYGZlt{\discretionary{}{\Wrappedafterbreak\char`\<}{\char`\<}}%
            \def\PYGZgt{\discretionary{\char`\>}{\Wrappedafterbreak}{\char`\>}}%
            \def\PYGZsh{\discretionary{}{\Wrappedafterbreak\char`\#}{\char`\#}}%
            \def\PYGZpc{\discretionary{}{\Wrappedafterbreak\char`\%}{\char`\%}}%
            \def\PYGZdl{\discretionary{}{\Wrappedafterbreak\char`\$}{\char`\$}}%
            \def\PYGZhy{\discretionary{\char`\-}{\Wrappedafterbreak}{\char`\-}}%
            \def\PYGZsq{\discretionary{}{\Wrappedafterbreak\textquotesingle}{\textquotesingle}}%
            \def\PYGZdq{\discretionary{}{\Wrappedafterbreak\char`\"}{\char`\"}}%
            \def\PYGZti{\discretionary{\char`\~}{\Wrappedafterbreak}{\char`\~}}%
        }
        % Some characters . , ; ? ! / are not pygmentized.
        % This macro makes them "active" and they will insert potential linebreaks
        \newcommand*\Wrappedbreaksatpunct {%
            \lccode`\~`\.\lowercase{\def~}{\discretionary{\hbox{\char`\.}}{\Wrappedafterbreak}{\hbox{\char`\.}}}%
            \lccode`\~`\,\lowercase{\def~}{\discretionary{\hbox{\char`\,}}{\Wrappedafterbreak}{\hbox{\char`\,}}}%
            \lccode`\~`\;\lowercase{\def~}{\discretionary{\hbox{\char`\;}}{\Wrappedafterbreak}{\hbox{\char`\;}}}%
            \lccode`\~`\:\lowercase{\def~}{\discretionary{\hbox{\char`\:}}{\Wrappedafterbreak}{\hbox{\char`\:}}}%
            \lccode`\~`\?\lowercase{\def~}{\discretionary{\hbox{\char`\?}}{\Wrappedafterbreak}{\hbox{\char`\?}}}%
            \lccode`\~`\!\lowercase{\def~}{\discretionary{\hbox{\char`\!}}{\Wrappedafterbreak}{\hbox{\char`\!}}}%
            \lccode`\~`\/\lowercase{\def~}{\discretionary{\hbox{\char`\/}}{\Wrappedafterbreak}{\hbox{\char`\/}}}%
            \catcode`\.\active
            \catcode`\,\active
            \catcode`\;\active
            \catcode`\:\active
            \catcode`\?\active
            \catcode`\!\active
            \catcode`\/\active
            \lccode`\~`\~
        }
    \makeatother

    \let\OriginalVerbatim=\Verbatim
    \makeatletter
    \renewcommand{\Verbatim}[1][1]{%
        %\parskip\z@skip
        \sbox\Wrappedcontinuationbox {\Wrappedcontinuationsymbol}%
        \sbox\Wrappedvisiblespacebox {\FV@SetupFont\Wrappedvisiblespace}%
        \def\FancyVerbFormatLine ##1{\hsize\linewidth
            \vtop{\raggedright\hyphenpenalty\z@\exhyphenpenalty\z@
                \doublehyphendemerits\z@\finalhyphendemerits\z@
                \strut ##1\strut}%
        }%
        % If the linebreak is at a space, the latter will be displayed as visible
        % space at end of first line, and a continuation symbol starts next line.
        % Stretch/shrink are however usually zero for typewriter font.
        \def\FV@Space {%
            \nobreak\hskip\z@ plus\fontdimen3\font minus\fontdimen4\font
            \discretionary{\copy\Wrappedvisiblespacebox}{\Wrappedafterbreak}
            {\kern\fontdimen2\font}%
        }%

        % Allow breaks at special characters using \PYG... macros.
        \Wrappedbreaksatspecials
        % Breaks at punctuation characters . , ; ? ! and / need catcode=\active
        \OriginalVerbatim[#1,codes*=\Wrappedbreaksatpunct]%
    }
    \makeatother

    % Exact colors from NB
    \definecolor{incolor}{HTML}{303F9F}
    \definecolor{outcolor}{HTML}{D84315}
    \definecolor{cellborder}{HTML}{CFCFCF}
    \definecolor{cellbackground}{HTML}{F7F7F7}

    % prompt
    \makeatletter
    \newcommand{\boxspacing}{\kern\kvtcb@left@rule\kern\kvtcb@boxsep}
    \makeatother
    \newcommand{\prompt}[4]{
        {\ttfamily\llap{{\color{#2}[#3]:\hspace{3pt}#4}}\vspace{-\baselineskip}}
    }
    

    
    % Prevent overflowing lines due to hard-to-break entities
    \sloppy
    % Setup hyperref package
    \hypersetup{
      breaklinks=true,  % so long urls are correctly broken across lines
      colorlinks=true,
      urlcolor=urlcolor,
      linkcolor=linkcolor,
      citecolor=citecolor,
      }
    % Slightly bigger margins than the latex defaults
    
    \geometry{verbose,tmargin=1in,bmargin=1in,lmargin=1in,rmargin=1in}
    
    

\begin{document}
    
    \maketitle
    
    

    
    \begin{tcolorbox}[breakable, size=fbox, boxrule=1pt, pad at break*=1mm,colback=cellbackground, colframe=cellborder]
\prompt{In}{incolor}{1}{\boxspacing}
\begin{Verbatim}[commandchars=\\\{\}]
\PY{k+kn}{import}\PY{+w}{ }\PY{n+nn}{marimo}\PY{+w}{ }\PY{k}{as}\PY{+w}{ }\PY{n+nn}{mo}
\PY{k+kn}{import}\PY{+w}{ }\PY{n+nn}{numpy}\PY{+w}{ }\PY{k}{as}\PY{+w}{ }\PY{n+nn}{np}
\end{Verbatim}
\end{tcolorbox}

    \section{Mandat no. 1 :
Probabilités}\label{mandat-no.-1-probabilituxe9s}

    \subsection{Partie 1: Roues de
pouvoirs}\label{partie-1-roues-de-pouvoirs}

    \subsubsection{i)}\label{i}

Le pouvoir est obtenu si toutes les roues s'arrêtent sur le même
pictogramme. Deux configurations sont testés: - 3 roues avec 8
pictogrammes chacune - 4 roues avec 5 pictogrammes chacune

C'est un cas classique de probabilités basé sur le principe fondamental
du dénombrement, avec des événements indépendants.

\paragraph{Premier cas : 3 roues, 8
pictogrammes}\label{premier-cas-3-roues-8-pictogrammes}

L'espace échantillonnal contient 8 pictogrammes. L'événement recherché
est d'obtenir le même pictogramme sur les 3 roues.

La probabilité d'obtenir un pictogramme spécifique sur les 3 roues est:
\[
\frac{1}{8} \times \frac{1}{8} \times \frac{1}{8} = \frac{1}{512} \approx 0.195\%
\]

Cependant, c'est un seul pictogramme et il y en a 8 possibles, la
probabilité d'obtenir un pouvoir (peu importe lequel) est donc en
fonction du nombre total de possibilités \(8^3\) et du nombre de
possibilités valides (\(8\)): \[
\frac{8}{512} \approx 1.56\%
\]

\paragraph{Deuxième cas : 4 roues, 5
pictogrammes}\label{deuxiuxe8me-cas-4-roues-5-pictogrammes}

Même chose qu'en haut: La probabilité d'obtenir un pictogramme
spécifique sur les 4 roues est: \[
\frac{1}{5} \times \frac{1}{5} \times \frac{1}{5} \times \frac{1}{5} = \frac{1}{625} = 0.16\%
\]

Pour les 5 possibles, la probabilité d'obtenir un pouvoir c'est: \[
\frac{5}{625} \approx 0.8\%
\]

\paragraph{Résultats}\label{ruxe9sultats}

La configuration avec 3 roues et 8 pictogrammes offre une probabilité
plus élevée d'obtenir un pouvoir.

    \subsubsection{ii)}\label{ii}

C'est le cas où un pouvoir est obtenue si toutes les roues affichent des
pictogrammes différents.

\paragraph{Permutations ou
combinaisons}\label{permutations-ou-combinaisons}

Pour déterminer le nombre de possibilités, il faut choisir entre
permutations ou combinaisons. Une combinaison ne tient pas compte de
l'ordre, tandis qu'une permutation en tient compte.

Les permutations sont utilisé puisque l'ordre est important: par
exemple, (A, B, C) et (C, A, B) sont deux résultats différents. \[
P(n,r)=\frac{n!}{(n-r)!}
\]

où \(n\) est le nombre de pictogrammes possibles et \(r\) le nombre de
roues.

\paragraph{Premier cas : 3 roues, 8
pictogrammes}\label{premier-cas-3-roues-8-pictogrammes}

\[
P(8,3)=\frac{8!}{(8-3)!}=\frac{8!}{5!}=336
\] Il y a donc 336 arrangements où les pictogrammes sont tous
différents.

Sachant que le nombre total d'arrangement est \(8^3=512\), La
probabilité d'obtenir uniquement des pictogrammes différents est donc:
\[
\frac{336}{512} \approx 65.6\%
\]

\paragraph{Deuxième cas : 4 roues, 5
pictogrammes}\label{deuxiuxe8me-cas-4-roues-5-pictogrammes}

\[
P(5,4)=\frac{5!}{(5-4)!}=\frac{5!}{1!}=120
\] Il y a donc 120 arrangements où les pictogrammes sont tous
différents.

Sachant que le nombre total d'arrangement est \(5^4=625\), La
probabilité d'obtenir uniquement des pictogrammes différents est donc:
\[
\frac{120}{625} \approx 19.2\%
\]

\paragraph{Résultats}\label{ruxe9sultats}

Il est beaucoup plus facile d'obtenir des pictogrammes tous différents
dans le premier cas que dans le deuxième.

    \subsubsection{iii)}\label{iii}

Dans ce cas, on obtient un pouvoir si toutes les roues affichent le même
pictogramme au moins 2 fois sur 5 essais. Les succès peuvent être
différents (par exemple : (AAA) et (BBB) sont valides).

On modélise la situation avec une variable aléatoire \(X\) représentant
le nombre de succès (Nombre de fois où toutes les roues affichent le
même pictogramme).

On utilise la loi binomiale, car : - le nombre d'essais est fixé
(\(n = 5\)) - chaque essai a deux issues (succès ou échec) - la
probabilité de succès est constante - les essais sont indépendants

La fonction de masse de la loi binomiale c'est: \[
P(X=k)=\binom{n}{k} p^k (1-p)^{n-k}
\]

où: \[
\binom{n}{k}=\frac{n!}{k!(n-k)!}
\]

On cherche: \[
P(X \ge 2)
\]

\paragraph{Probabilitées}\label{probabilituxe9es}

\subparagraph{Premier cas: (3 roues, 8
pictogrammes)}\label{premier-cas-3-roues-8-pictogrammes}

La probabilité de succès d'avoir que des pictogrammes identiques à déjà
été calculé en haut. C'est \(1.56\%\).

On calcul les probabilitées d'en avoir moins que 2. C'est plus court et
le résultat n'a qu'a être inversé. \[
P(X=0)=\frac{5!}{0!(5-0)!}0.0156^0(1-0.0156)^{5-0}=0.92439593103
\] \[
P(X=1)=\frac{5!}{1!(5-1)!}0.0156^1(1-0.0156)^{5-1}=0.07324551261
\]

l'inverse, est donc: \[
P(X \ge 2)=1 - (P(X=0)+P(X=1))
\] \[
P(X \ge 2)=1 - (0.9244 + 0.0732) \approx 0.0024
\] Soit environ \(0.24\%\)

\subparagraph{Deuxième cas: (4 roues, 5
pictogrammes)}\label{deuxiuxe8me-cas-4-roues-5-pictogrammes}

La probabilité de succès d'avoir que des pictogrammes identiques à déjà
été calculé en haut. C'est \(0.8\%\). On calcul les probabilitées d'en
avoir moins que 2. C'est plus court et le résultat n'a qu'a être
inversé. \[
P(X=0)=\frac{5!}{0!(5-0)!}0.008^0(1-0.008)^{5-0}=0.96063490044
\] \[
P(X=1)=\frac{5!}{1!(5-1)!}0.008^1(1-0.008)^{5-1}=0.03873527824
\]

l'inverse, est donc: \[
P(X \ge 2)=1 - (P(X=0)+P(X=1))
\] \[
P(X \ge 2)=1 - (0.9606 + 0.0387) \approx 0.0006
\] Soit environ: \(0.06\%\)

\paragraph{Variance}\label{variance}

Pour une loi binomiale, la variance est donnée par: \[
\mathrm{Var}(X)=n \cdot p (1-p)
\]

\subparagraph{Premier cas:}\label{premier-cas}

\[
\mathrm{Var}(X)=5 \cdot 0.0156 \cdot (1-0.0156) \approx 0.0768
\]

\subparagraph{Deuxième cas:}\label{deuxiuxe8me-cas}

\[
\mathrm{Var}(X)=5 \cdot 0.008 \cdot (1-0.008) \approx 0.0397
\]

\paragraph{Espérance}\label{espuxe9rance}

L'espérance de la loi binomiale c'est: \[
E[X]=n \cdot p
\]

\subparagraph{Premier cas:}\label{premier-cas-1}

\[
E[X]=5 \cdot 0.0156 \approx 0.078
\]

\subparagraph{Deuxième cas:}\label{deuxiuxe8me-cas-1}

\[
E[X]=5 \cdot 0.008 \approx 0.04
\]

    \subsection{Partie 2: Probabilitées de la
cible}\label{partie-2-probabilituxe9es-de-la-cible}

    \subsubsection{\texorpdfstring{Covariance entre \(X\) et
\(Y\)}{Covariance entre X et Y}}\label{covariance-entre-x-et-y}

Intuitivement, il n'y a pas de relation entre \(X\) et \(Y\). Un
déplacement selon l'axe \(X\) n'affecte pas la position selon \(Y\), et
inversement. On suppose donc que : \[
\mathrm{Cov}(X,Y) = 0
\]

\paragraph{\texorpdfstring{Corrélation entre \(X\) et
\(Y\)}{Corrélation entre X et Y}}\label{corruxe9lation-entre-x-et-y}

La corrélation est donnée par : \[
\rho_{XY}=\frac{\mathrm{Cov}(X,Y)}{\sigma_X \sigma_Y}
\] Puisque la covariance est nulle : \[
\rho_{XY} = 0
\] Dans le cas d'une loi normale bivariée, une corrélation nulle
implique que les variables sont indépendantes. Donc, \(X\) et \(Y\) sont
indépendantes.

\paragraph{Fonction de densité de
probabilité}\label{fonction-de-densituxe9-de-probabilituxe9}

La fonction de densité de probabilité conjointe pour deux variables
indépendantes \(X\) et \(Y\) (loi normale bivariée) s'écrit: \[
f_{X,Y}(x,y|z) = f_X(x) \cdot f_Y(y)
\]

La PDF d'une loi normale univariée est : \[
f(x) = \frac{1}{\sigma \sqrt{2\pi}} \, e^{-\frac{(x-\mu)^2}{2\sigma^2}}
\]

Donc, pour la distribution conjointe :

\[
f_{X,Y}(x,y|z) =
\frac{1}{\sigma_X \sqrt{2\pi}} \, e^{-\frac{(x-\mu_X)^2}{2\sigma_X^2}}
\times
\frac{1}{\sigma_Y \sqrt{2\pi}} \, e^{-\frac{(y-\mu_Y)^2}{2\sigma_Y^2}}
\]

Comme \(\mu_X = \mu_Y = 0\), ça se simplifie:

\[
f_{X,Y}(x,y|z) =
\frac{1}{\sigma_X \sqrt{2\pi}} \, e^{-\frac{x^2}{2\sigma_X^2}}
\times
\frac{1}{\sigma_Y \sqrt{2\pi}} \, e^{-\frac{y^2}{2\sigma_Y^2}}
\]

À noter : \(\sigma_Y\) dépend de la variable \(z\) (distance par rapport
à la cible).

\paragraph{Fonction de densité de probabilité
marginale}\label{fonction-de-densituxe9-de-probabilituxe9-marginale}

Comme \(X\) et \(Y\) sont indépendant, il est possible d'obtenir les PDF
marginales simplement en ``coupant'' la PDF conjointe. \#\#\#\#\# Pour
\(X\): \[
f_X(x) =
\frac{1}{\sigma_X \sqrt{2\pi}} \, e^{-\frac{x^2}{2\sigma_X^2}}
\] Avec \(\sigma_X = 0.1\): \[
f_X(x) =
\frac{1}{0.1 \sqrt{2\pi}} \, e^{-\frac{x^2}{0.02}}
\] Ce qui donne: \[
f_X(x) \approx 4 \, e^{-\frac{x^2}{0.02}}
\]

\subparagraph{\texorpdfstring{Pour \(Y\):}{Pour Y:}}\label{pour-y}

Pour \(Y\), on garde le \(\sigma\) puisqu'il peut changé en fonction de
\(z\): \[
f_Y(y) =
\frac{1}{\sigma_Y \sqrt{2\pi}} \, e^{-\frac{y^2}{2\sigma_Y^2}}
\] \[
f_Y(y|z) =
\frac{1}{\sigma_Y(z) \sqrt{2\pi}} \, e^{-\frac{y^2}{2\sigma_Y(z)^2}}
\]

\paragraph{Début de la probabilité d'ouvrir la porte à moins d'un
mètre}\label{duxe9but-de-la-probabilituxe9-douvrir-la-porte-uxe0-moins-dun-muxe8tre}

Pour réussir, il faut que les coordonnées du lancé tombent à l'intérieur
d'un cercle de rayon de \(\le 0.1\). La condition est simple grâce au
théorème de Pythagore: \[
X^2 + Y^2 \le R^2
\] \[
X^2 + Y^2 \le 0.1^2
\]

Comme on connaît \(\sigma_Y\) pour cette distance, la PDF de \(Y\) peut
être simplifiée : \[
f_Y(y|z) = \frac{1}{\sigma_Y(z) \sqrt{2\pi}} \, e^{-\frac{y^2}{2\sigma_Y(z)^2}}
\] \[
f_Y(y) = \frac{1}{0.05 \sqrt{2\pi}} \, e^{-\frac{y^2}{2\cdot 0.05^2}} \approx 8 \, e^{-\frac{y^2}{0.005}}
\]

\subparagraph{Pourquoi on ne fait pas l'intégrale à la
main}\label{pourquoi-on-ne-fait-pas-lintuxe9grale-uxe0-la-main}

La probabilité qu'une variable continue tombe dans une région donnée se
calcule avec son \textbf{CDF}, ça implique une \textbf{intégrale}. Dans
notre cas, il s'agit d'une double intégrale quand même assez complex:

\[
P(X,Y) = \int_{-\infty}^{\infty} \int_{-\infty}^{\infty} f_{X,Y}(x,y) \, dx \, dy
\]

\(X\) et \(Y\) étant indépendant,la PDF conjointe se factorise: \[
F_X(x) = \int_{-\infty}^{\infty} f_X(x) \, dx, \quad
F_Y(y) = \int_{-\infty}^{\infty} f_Y(y) \, dy
\]

Les intégrales restent lourds à faire à la main. Donc Monte-Carlo est
utilisé, vue que la condition de réussite est très simple, les
distributions, variances et moyennes sont connues.

\subparagraph{Estimation par simulation
Monte-Carlo}\label{estimation-par-simulation-monte-carlo}

Avec Monte-Carlo, on peut estimer la probabilité beaucoup plus
facilement. La probabilité devient simplement : \[
P \approx \frac{\text{Nombre de succès}}{N}
\]

Plus le nombre de tirages \(N\) est grand, plus l'estimation est
précise. De plus, NumPy à déjà une fonction pour généré des valeurs
aléatoires dans une distribution normale.

Hors le processus est simple. N points aléatoire distribués normalements
avec une moyenne de \(0\) sont généré et avec la variance connue.
Ensuites, ils sont comparé avec la condition de réussite, qui est que la
distance entre le centre du cercle et la coordonée soit égal ou moindre
que \(0.1\). Le taux de succès sur le nombre total d'essaie donne la
probabilité de réussite.

    \begin{tcolorbox}[breakable, size=fbox, boxrule=1pt, pad at break*=1mm,colback=cellbackground, colframe=cellborder]
\prompt{In}{incolor}{2}{\boxspacing}
\begin{Verbatim}[commandchars=\\\{\}]
\PY{k}{def}\PY{+w}{ }\PY{n+nf}{random\PYZus{}value}\PY{p}{(}\PY{n}{variance}\PY{p}{,} \PY{n}{N}\PY{p}{)}\PY{p}{:}
    \PY{k}{return} \PY{n}{np}\PY{o}{.}\PY{n}{random}\PY{o}{.}\PY{n}{normal}\PY{p}{(}\PY{l+m+mi}{0}\PY{p}{,} \PY{n}{variance}\PY{p}{,} \PY{n}{N}\PY{p}{)}


\PY{k}{def}\PY{+w}{ }\PY{n+nf}{is\PYZus{}in\PYZus{}circle}\PY{p}{(}\PY{n}{x}\PY{p}{,} \PY{n}{y}\PY{p}{,} \PY{n}{radius}\PY{p}{)}\PY{p}{:}
    \PY{k}{return} \PY{p}{(}\PY{n}{x}\PY{o}{*}\PY{o}{*}\PY{l+m+mi}{2} \PY{o}{+} \PY{n}{y}\PY{o}{*}\PY{o}{*}\PY{l+m+mi}{2}\PY{p}{)} \PY{o}{\PYZlt{}}\PY{o}{=} \PY{n}{radius}\PY{o}{*}\PY{o}{*}\PY{l+m+mi}{2}
\end{Verbatim}
\end{tcolorbox}

    \begin{tcolorbox}[breakable, size=fbox, boxrule=1pt, pad at break*=1mm,colback=cellbackground, colframe=cellborder]
\prompt{In}{incolor}{3}{\boxspacing}
\begin{Verbatim}[commandchars=\\\{\}]
\PY{n}{N} \PY{o}{=} \PY{l+m+mi}{20000000}

\PY{n}{X} \PY{o}{=} \PY{n}{random\PYZus{}value}\PY{p}{(}\PY{l+m+mf}{0.1}\PY{p}{,} \PY{n}{N}\PY{p}{)}
\PY{n}{Y} \PY{o}{=} \PY{n}{random\PYZus{}value}\PY{p}{(}\PY{l+m+mf}{0.05}\PY{p}{,} \PY{n}{N}\PY{p}{)}
\PY{n}{results} \PY{o}{=} \PY{n}{is\PYZus{}in\PYZus{}circle}\PY{p}{(}\PY{n}{X}\PY{p}{,} \PY{n}{Y}\PY{p}{,} \PY{l+m+mf}{0.1}\PY{p}{)}
\PY{n}{success} \PY{o}{=} \PY{n}{np}\PY{o}{.}\PY{n}{sum}\PY{p}{(}\PY{n}{results}\PY{p}{)}
\PY{n}{probability} \PY{o}{=} \PY{n}{success} \PY{o}{/} \PY{n}{N}

\PY{n+nb}{print}\PY{p}{(}\PY{l+s+s2}{\PYZdq{}}\PY{l+s+s2}{z\PYZlt{}1: }\PY{l+s+s2}{\PYZdq{}}\PY{p}{,} \PY{n}{probability}\PY{p}{)}

\PY{n}{X} \PY{o}{=} \PY{n}{random\PYZus{}value}\PY{p}{(}\PY{l+m+mf}{0.1}\PY{p}{,} \PY{n}{N}\PY{p}{)}
\PY{n}{Y} \PY{o}{=} \PY{n}{random\PYZus{}value}\PY{p}{(}\PY{l+m+mf}{0.4}\PY{p}{,} \PY{n}{N}\PY{p}{)}
\PY{n}{results} \PY{o}{=} \PY{n}{is\PYZus{}in\PYZus{}circle}\PY{p}{(}\PY{n}{X}\PY{p}{,} \PY{n}{Y}\PY{p}{,} \PY{l+m+mf}{0.1}\PY{p}{)}
\PY{n}{success} \PY{o}{=} \PY{n}{np}\PY{o}{.}\PY{n}{sum}\PY{p}{(}\PY{n}{results}\PY{p}{)}
\PY{n}{probability} \PY{o}{=} \PY{n}{success} \PY{o}{/} \PY{n}{N}

\PY{n+nb}{print}\PY{p}{(}\PY{l+s+s2}{\PYZdq{}}\PY{l+s+s2}{z\PYZgt{}10: }\PY{l+s+s2}{\PYZdq{}}\PY{p}{,} \PY{n}{probability}\PY{p}{)}
\end{Verbatim}
\end{tcolorbox}

    \begin{Verbatim}[commandchars=\\\{\}]
z<1:  0.59012135
z>10:  0.1103074
    \end{Verbatim}

    \begin{tcolorbox}[breakable, size=fbox, boxrule=1pt, pad at break*=1mm,colback=cellbackground, colframe=cellborder]
\prompt{In}{incolor}{ }{\boxspacing}
\begin{Verbatim}[commandchars=\\\{\}]

\end{Verbatim}
\end{tcolorbox}


    % Add a bibliography block to the postdoc
    
    
    
\end{document}
